\chapter{Metode Gauss}

\Sage{} dapat menyelesaikan sistem linear dengan beberapa cara.
Program ini dapat menggunakan pemecah sistem umum, dan
juga dapat menggunakan pemecah yang dikhususkan untuk sistem linear.
Kita akan melihat keduanya. 



\section{Sistem persamaan}
Untuk memasukkan suatu sistem persamaan, pertama-tama Anda harus memasukkan
persamaan-persamaan tunggal.
Jadi, Anda harus memulai dengan variabel.
Kita telah melihat satu jenis variabel ketika memberikan perintah seperti berikut.
\begin{sagecommandline}
sage: x = 3
sage: 7*x
\end{sagecommandline}
Seperti dibahas dalam bab sebelumnya,
di sini $x$ adalah nama suatu lokasi dalam memori komputer
yang kita gunakan untuk menyimpan dan mengambil nilai.

Variabel dalam persamaan berbeda; dalam persamaan
$C=2\pi\cdot r$, kedua variabel tidak mempunyai nilai tetap dan
tidak terikat pada lokasi memori. 
Untuk memperlihatkan perbedaannya, ketik variabel yang belum diberi nilai,
misalnya \inlinecode{y}, lalu tekan \keyboardkey{Enter}.
Anda akan memperoleh pengecualian yang baris terakhirnya berbunyi
\inlinecode{NameError: name 'y' is not defined}.

Untuk menggunakan $y$ sebagai variabel simbolis, pertama-tama Anda harus
mendeklarasikannya kepada sistem.
\begin{sagecommandline}
sage: var('y')
sage: y
sage: 2*y
\end{sagecommandline}
Sebelumnya, \Sage{} memperlakukan $x$ sebagai lokasi untuk menyimpan
nilai sehingga mengevaluasi \inlinecode{7*x} (dan memperoleh nilai $21$).
Namun, \Sage{} tidak mengevaluasi \inlinecode{2*y} karena
Anda telah memberitahunya bahwa $y$ adalah variabel simbolis.

Dengan demikian, suatu sistem persamaan merupakan sebuah daftar.
\begin{sagecommandline}
sage: var('x,y,z')                                  
sage: eqns = [x-y+2*z == 4, 2*x+2*y == 12, x-4*z==5]
\end{sagecommandline}
Ada dua hal yang perlu diperhatikan di sini:~dalam persamaan, Anda
harus menulis dua tanda sama dengan \inlinecode{==}, bukan
operator penugasan \inlinecode{=},
dan Anda harus menulis \inlinecode{2*x},
bukan \inlinecode{2x}.
Kedua jenis kesalahan tersebut akan memicu pesan galat yang berbunyi
\inlinecode{SyntaxError: invalid syntax}.

Salah satu cara menyelesaikan sistem ini adalah dengan pemecah umum milik \Sage{}.
\begin{sagecommandline}
sage: var('x,y,z')                                  
sage: eqns = [x-y+2*z == 4, 2*x+2*y == 12, x-4*z==5]
sage: solve(eqns, x, y, z)                            
\end{sagecommandline}
Hasilnya cukup baik, tetapi pemecah ini cukup canggih dan dapat memberikan
jawaban yang bukan sekadar bilangan.
Selanjutnya, kita menaruh parameter~\inlinecode{a} di ruas kanan. 
\Sage{}
menyatakan penyelesaian bagi $x$, $y$, dan~$z$ dalam parameter tersebut.
\begin{sagecommandline}
sage: var('x,y,z,a')                                
sage: eqns = [x-y+2*z == a, 2*x+2*y == 12, x-4*z==5]
sage: solve(eqns, x, y, z) 
\end{sagecommandline}



\subsection{Matriks}
Rutin \inlinecode{solve} bersifat umum, tetapi
untuk kasus khusus sistem linear tersedia alat yang lebih baik.
Kita akan membahas matriks dan vektor lebih lanjut dalam bab-bab berikutnya,
tetapi secara ringkas:~matriks \Sage{} adalah daftar baris, dengan
setiap baris berupa daftar bilangan. 
Kita mendeklarasikan jenis bilangannya di awal.
Dalam buku ini, umumnya kita menganggap entri-entri sebagai bilangan real,
tetapi contoh dan latihan menggunakan bilangan rasional yang
lebih mudah dibaca.
\Sage{} menggunakan \Sagecmd{QQ} untuk himpunan bilangan rasional.
\begin{sagecommandline}
sage: M = matrix(QQ, [[1, 2, 3, 4], [5, 6, 7, 8], [9, 10, 11, 12]])
sage: M
sage: M.nrows()
sage: M.ncols()
sage: M[1,2]
\end{sagecommandline}
Seperti daftar \python{}, daftar \Sage{} diindeks mulai dari nol.  
Jadi, \inlinecode{M[1,2]} meminta
entri pada baris kedua dan kolom ketiga. 

Masukkan vektor dengan cara yang hampir sama.
\begin{sagecommandline}
sage: v = vector(QQ, [2/3, -1/3, 1/2])
sage: v
sage: v[1]
\end{sagecommandline}
\Sage{} tidak mempermasalahkan perbedaan antara vektor baris dan
vektor kolom.
Di sini, \inlinecode{v} dicetak dengan tanda kurung biasa, berbeda dengan
matriks yang dicetak dengan tanda kurung siku.
Jika Anda perlu memastikan apakah yang dimiliki adalah baris atau kolom,
ubahlah vektor tersebut menjadi matriks.
\begin{sagecommandline}
sage: v = vector(QQ, [1,3,5])
sage: v.row()
sage: v.column()
\end{sagecommandline}

Buku ini sering menyelesaikan sistem linear dengan memperluas matriks
menggunakan suatu vektor.
\begin{sagecommandline}
sage: M = matrix(QQ, [[1, 2, 3], [4, 5, 6], [7 ,8, 9]])
sage: v = vector(QQ, [2/3, -1/3, 1/2])
sage: M_prime = M.augment(v)
sage: M_prime
\end{sagecommandline}
Argumen opsional membuat
\Sage{} menampilkan pemisah antara kedua bagian
$M^\prime$.
\begin{sagecommandline}
sage: M = matrix(QQ, [[1, 2, 3], [4, 5, 6], [7 ,8, 9]])
sage: v = vector(QQ, [2/3, -1/3, 1/2])
sage: M_prime = M.augment(v, subdivide=True)
sage: M_prime                              
\end{sagecommandline}



\subsection{Operasi baris}
%% Computers are a big help with jobs that are tedious and error-prone.
%% Row operations are both.
Kita dapat meminta \Sage{} mengerjakan aritmetika dalam Metode Gauss.
%% First enter an example matrix.
\begin{sagecommandline}
sage: M = matrix(QQ, [[0, 2, 1], [2, 0, 4], [2 ,-1/2, 3]])
sage: v = vector(QQ, [2, 1, -1/2])                        
sage: M_prime = M.augment(v, subdivide=True)              
sage: M_prime                                             
\end{sagecommandline}
Tukarkan dua baris teratas. 
Ingat bahwa indeks daftar dimulai dari nol,
sehingga baris teratas adalah baris~$0$.
\begin{sagecommandline}
sage: M = matrix(QQ, [[0, 2, 1], [2, 0, 4], [2 ,-1/2, 3]])
sage: v = vector(QQ, [2, 1, -1/2])                        
sage: M_prime = M.augment(v, subdivide=True)              
sage: M_prime.swap_rows(0,1)
sage: M_prime
\end{sagecommandline}
Selanjutnya, skalakan ulang baris teratas dengan mengalikan entri-entrinya
dengan~$1/2$.
\begin{sagecommandline}
sage: M_prime.rescale_row(0, 1/2)
sage: M_prime
\end{sagecommandline}
Ganti baris terbawah dengan
hasil penjumlahannya dengan
$-2$ kali baris teratas.
\begin{sagecommandline}
sage: M_prime.add_multiple_of_row(2,0,-2)
sage: M_prime
\end{sagecommandline}
% One more row combination operation gives echelon form.
Terakhir, ganti baris terbawah dengan hasil penjumlahannya dengan
$1/4$ kali baris tengah.
\begin{sagecommandline}
sage: M_prime.add_multiple_of_row(2,1,1/4)
sage: M_prime                             
\end{sagecommandline}
Sekarang, substitusi balik secara manual akan memberikan penyelesaiannya,
atau kita dapat menggunakan \inlinecode{solve} milik \Sage{}.
\begin{sagecommandline}
sage: var('x,y,z')
sage: eqns=[-3/4*z == -1, 2*y+z == 2, x+2*z == 1/2]
sage: solve(eqns, x, y, z)
\end{sagecommandline}

Operasi-operasi di atas
bekerja secara langsung pada objeknya.
Artinya, operasi tersebut mengubah matriks~$M^\prime$.
\Sage{} mempunyai perintah terkait yang membiarkan matriks awal tidak berubah,
tetapi mengembalikan matriks yang telah diubah.
\begin{sagecommandline}
sage: M = matrix(QQ, [[0, 1, -1], [1, -2, 0], [0 ,-1, 4]])
sage: v = vector(QQ, [0, 1, -2])
sage: M_prime = M.augment(v, subdivide=True) 
sage: M_prime
sage: N = M_prime.with_swapped_rows(0,1)
sage: M_prime
sage: N      
\end{sagecommandline}
Di sini, rutin tersebut tidak mengubah $M^\prime$, sedangkan $N$ adalah
matriks hasil perubahan yang dikembalikan.
Dua rutin lain dari jenis ini adalah \inlinecode{with_rescaled_row}
dan \inlinecode{with_added_multiple_of_row}.




\subsection{Sistem nonsingular dan singular}
Langkah-langkah reduksi matriks yang kita jalani bersifat mekanis.
Artinya, kita semestinya dapat mengotomatiskan proses dari suatu matriks
ke bentuk eselon baris tereduksinya.
\begin{sagecommandline}
sage: M = matrix(QQ, [[1/2, 1, -1], [1, -2, 0], [2 ,-1, 1]])
sage: v = vector(QQ, [0, 1, -2])
sage: M_prime = M.augment(v, subdivide=True) 
sage: M_prime.rref()
\end{sagecommandline}
Jadi, kita telah memperoleh jawaban tanpa \inlinecode{solve} milik \Sage{}.
Keuntungan menggunakan metode yang dikhususkan
untuk sistem linear ialah metode tersebut lebih cepat dan juga dapat lebih jarang
mengalami masalah aritmetika titik mengambang.

Dalam contoh tersebut, $M$ nonsingular karena berbentuk
persegi dan, setelah diubah ke bentuk eselon, setiap
kolom bersesuaian dengan variabel utama.
Contoh berikutnya mempunyai
matriks persegi yang singular. 
% Consequently, in the  
% echelon form system there are columns on the left 
% that do not have a leading variable. 
\begin{sagecommandline}
sage: M = matrix(QQ, [[1, 1, 1, 1], [1, 2, 3, 4], [2, 3, 4, 5], [0, 1, 2 ,3]]) 
sage: v = vector(QQ, [0, 1, 1, 1]) 
sage: M_prime = M.augment(v, subdivide=True)
sage: M_prime
sage: M_prime.rref()
\end{sagecommandline}
Ingat bahwa jika suatu sistem linear mempunyai matriks koefisien persegi
yang singular, terdapat dua kemungkinan.
Kemungkinan pertama tampak di atas:~dalam bentuk eselon,
setiap baris yang seluruh entri di sebelah kirinya nol
juga mempunyai entri nol di sebelah kanan.
Sistem ini mempunyai tak berhingga banyak penyelesaian.

Sebaliknya, dengan matriks awal yang sama,
contoh di bawah mempunyai baris yang entri-entri di sebelah kirinya nol,
tetapi entri di sebelah kanannya tak nol, sehingga tidak mempunyai penyelesaian.
\begin{sagecommandline} 
sage: v = vector(QQ, [0, 1, 2, 1])             
sage: M_prime = M.augment(v, subdivide=True)
sage: M_prime.rref()                        
\end{sagecommandline}

Perbedaan antara kedua kemungkinan tersebut
berkaitan dengan hubungan di antara
baris-baris~$M$ dan juga hubungan di antara entri-entri vektor.
Dalam matriks itu,
baris ketiga merupakan jumlah dua baris pertama, sedangkan baris keempat
merupakan selisih baris kedua dan baris pertama.
Untuk vektornya, pada kasus pertama entri-entri vektor mempunyai hubungan
yang sama\Dash entri ketiga vektor merupakan jumlah dua entri pertama,
dan entri keempat merupakan selisih entri kedua dan entri pertama\Dash
sedangkan pada kasus kedua, entri-entri vektor tidak mempunyai hubungan itu.

Cara mudah untuk memastikan bahwa baris nol pada matriks
di sebelah kiri berpasangan dengan entri nol
pada vektor di sebelah kanan adalah membuat semua entri vektor bernilai nol;
dengan kata lain, meninjau sistem homogen yang bersesuaian.
\begin{sagecommandline}
sage: v = zero_vector(QQ, 4)
sage: v
sage: M = matrix(QQ, [[1, 1, 1, 1], [1, 2, 3, 4], [2, 3, 4, 5], [0, 1, 2 ,3]]) 
sage: M_prime = M.augment(v, subdivide=True)
sage: M_prime
sage: M_prime.rref()
\end{sagecommandline}

Anda dapat memperoleh indeks kolom-kolom yang mempunyai entri utama dengan
metode \inlinecode{pivots}.
\begin{sagecommandline}
sage: M = matrix(QQ, [[1, 1, 1, 1], [1, 2, 3, 4], [2, 3, 4, 5], [0, 1, 2 ,3]]) 
sage: v = vector(QQ, [0, 1, 1, 1])
sage: M_prime = M.augment(v, subdivide=True)
sage: M_prime.rref()
sage: M_prime.pivots()         
sage: M_prime.nonpivots()         
\end{sagecommandline}





\subsection{Parametrisasi}
Di atas, kita menggunakan rutin umum \inlinecode{solve} milik \Sage{}.
Rutin itu mungkin bukan cara terbaik untuk mencari penyelesaian
suatu sistem linear karena alat khusus kemungkinan lebih cepat dan
dapat lebih akurat. 
Namun, \inlinecode{solve} berguna untuk memberikan himpunan penyelesaian
sistem yang mempunyai tak berhingga banyak penyelesaian melalui parametrisasi.

Sebagai ilustrasi, di bawah ini kita menyusun
sistem dengan tak berhingga banyak penyelesaian.
Kita memulai dengan matriks koefisien
yang jumlah dua baris teratasnya sama dengan baris terbawah,
lalu menggandengkan vektor yang mempunyai hubungan baris yang sama. 
Selanjutnya, kita mencari bentuk eselon baris tereduksi.
Kemudian, kita menyalin hasilnya menjadi sistem persamaan dan menerapkan
\inlinecode{solve} dengan dua cara berbeda.
\begin{sagecommandline}
sage: M = matrix(QQ, [[1, 1, 1], [1, 2, 3], [2 , 3, 4]])    
sage: v = vector(QQ, [1, 0, 1])                            
sage: M_prime = M.augment(v, subdivide=True)               
sage: M_prime.rref()
sage: var('x,y,z')          
sage: eqns = [x-z == 2, y+2*z == -1]
sage: solve(eqns, x, y)   
sage: solve(eqns, x, y, z)                                 
\end{sagecommandline}
Pemanggilan pertama dari kedua pemanggilan \inlinecode{solve} tersebut
meminta penyelesaian hanya untuk $x$ dan~$y$, sehingga \Sage{} memberikan
penyelesaian dalam parameter~$z$.
Pada pemanggilan kedua, \Sage{} menghasilkan parameternya sendiri.   




%========================================
\section{Otomatisasi}

Sebagai penutup, kita memberikan kode \python{} untuk dua fungsi
yang meniru langkah-langkah yang dilakukan seseorang ketika
mengerjakan Metode Gauss secara manual.

Berkas sumber skrip tersebut terdapat di bawah. 
Arti ``secara manual'' di sini ialah bahwa skrip tersebut
mengasumsikan matriks-matriksnya kecil dan mempunyai entri bilangan rasional.
Dengan demikian, kita tidak perlu mencemaskan masalah aritmetika titik mengambang.
Singkatnya, ini adalah contoh sederhana yang
cocok untuk pekerjaan rumah, tetapi belum siap untuk aplikasi nyata.


\subsection{Memuat dan menjalankan}
Pertama, berikut beberapa contoh pemanggilan.
Jalankan \Sage{} di direktori yang memuat berkas \path{gauss_method.sage}.
% \begin{sagecommandline}
% sage: load("gauss_method.sage")
% sage: M = matrix(QQ, [[1/2, 1, 4], [2, 4, -1], [1, 2, 0]])          
% sage: v = vector(QQ, [-2, 5, 4])
% sage: M_prime = M.augment(v, subdivide=True)  
% sage: gauss_method(M_prime)
% \end{sagecommandline}
% gauss_method.sage uses print which is not captured
\begin{lstlisting}
sage: load("gauss_method.sage")                                                 
sage: M = matrix(QQ, [[1/2, 1, 4], [2, 4, -1], [1, 2, 0]])                      
sage: v = vector(QQ, [-2, 5, 4])                                                
sage: M_prime = M.augment(v, subdivide=True)                                    
sage: gauss_method(M_prime)                                                     
[1/2   1   4| -2]
[  2   4  -1|  5]
[  1   2   0|  4]
 ambil -4 kali baris 1 ditambah baris 2
 ambil -2 kali baris 1 ditambah baris 3
[1/2   1   4| -2]
[  0   0 -17| 13]
[  0   0  -8|  8]
 ambil -8/17 kali baris 2 ditambah baris 3
[  1/2     1     4|   -2]
[    0     0   -17|   13]
[    0     0     0|32/17]
\end{lstlisting}

Selain \inlinecode{gauss_method}\!, berkas tersebut juga memuat fungsi
\inlinecode{gauss_jordan}
untuk melanjutkan proses hingga mencapai bentuk eselon baris tereduksi.
\begin{lstlisting}
sage: load("gauss_method.sage")                                                 
sage: M = matrix(QQ, [[1/2, 1, 4], [2, 4, -1], [1, 2, 0]])                      
sage: v = vector(QQ, [-2, 5, 4])                                                
sage: M_prime = M.augment(v, subdivide=True)                                    
sage: gauss_jordan(M_prime)                                                     
[1/2   1   4| -2]
[  2   4  -1|  5]
[  1   2   0|  4]
 ambil -4 kali baris 1 ditambah baris 2
 ambil -2 kali baris 1 ditambah baris 3
[1/2   1   4| -2]
[  0   0 -17| 13]
[  0   0  -8|  8]
 ambil -8/17 kali baris 2 ditambah baris 3
[  1/2     1     4|   -2]
[    0     0   -17|   13]
[    0     0     0|32/17]
 ambil 2 kali baris 1
 ambil -1/17 kali baris 2
 ambil 17/32 kali baris 3
[     1      2      8|    -4]
[     0      0      1|-13/17]
[     0      0      0|     1]
 ambil 4 kali baris 3 ditambah baris 1
 ambil 13/17 kali baris 3 ditambah baris 2
[1 2 8|0]
[0 0 1|0]
[0 0 0|1]
 ambil -8 kali baris 2 ditambah baris 1
[1 2 0|0]
[0 0 1|0]
[0 0 0|1]
\end{lstlisting}



\subsection{Sumber \protect\path{gauss_method.sage}}

Berikut adalah implementasi sederhana Metode Gauss dan
reduksi Gauss--Jordan.

\lstinputlisting{gauss_method.sage}
\endinput


TODO:
